% GENERATED FILE. Edit canonical lessons, then run npm run generate:books.
% canonical-source-hash: fnv1a64:7b5d5f0bad62e79d
% canonical-lessons: ZH-C280-yijian, ZH-W280-yi, ZH-C280-juede, ZH-C280-haoxiang, ZH-C280-zheme, ZH-C280-practice

\chapter{What I Think}
\label{ch:zh-what-i-think}
\hlchaptermodality{280}

\begin{chapteropening}
\textbf{By the end of this chapter:} \emph{I can say what I think, and ask what you think.}

Saying what you think with \zh{我觉得} and a whole sentence, a suggestion or an objection, as if, and so much, built on one new written character: \zh{意}.
\end{chapteropening}

\section[\texorpdfstring{\zh{意见} (yìjiàn)}{意见 (yìjiàn)}]{\zh{意见} (yìjiàn) — a suggestion; an objection}

\label{lesson:ZH-C280-yijian}



\begin{quote}
A new chapter, and it opens on what people put forward when they disagree.

\begin{itemize}
\raggedright
  \item \emph{Recall:} say \zh{他突然来了} — \emph{tā tūrán lái le} — \textbf{R1}, the chapter you have finished
  \item \emph{Recall:} say \emph{kànfǎ}, a view — \textbf{R3}, durable
  \item \emph{Recall:} say \emph{yǒuyìsi}, interesting — \textbf{R3}
\end{itemize}
\end{quote}

\subsection*{You'll want to know: \zh{意见}}
\textbf{\zh{意见}} — \emph{yìjiàn} — \textquotedblleft{}a suggestion; an objection\textquotedblright{}.

\begin{quote}
\textbf{\zh{你有什么意见}} \emph{nǐ yǒu shénme yìjiàn} — What do you suggest? Or: what's your objection?
\end{quote}

\begin{quote}
\textbf{\zh{我没有意见}} \emph{wǒ méiyǒu yìjiàn} — I've no objection; fine by me.
\end{quote}

The next lesson puts the first character in your hand. It is the yì of \emph{yǒuyìsi}.

\begin{cousinweb}
\textbf{\zh{意见}} (\emph{yìjiàn}) is built from \zh{意} 'intention' and \zh{见} 'to see', both inherited from Middle Chinese, \emph{'iH} and \emph{kenH} (Guangyun, in Baxter's transcription). The dictionary glosses the first 'will, intention' and the second 'to look'; 'a suggestion' or 'an objection' for the pair is the modern word's own sense.
\end{cousinweb}

\begin{grammarlens}[title={What you've built}]
\emph{kànfǎ} is your view of a matter; \zh{意见} is what you put forward about it, and often it is a point against.

It is a noun: you have one, \zh{有意见}, or you have none, \zh{没有意见}.
\end{grammarlens}

\subsection*{Your turn}
\begin{itemize}
\raggedright
  \item \emph{Say it:} \emph{yìjiàn}
  \item \emph{Say it:} \emph{yìjiàn}, once more
  \item \emph{Say it:} \zh{你有什么意见} — \emph{nǐ yǒu shénme yìjiàn} — What do you suggest?
  \item \emph{Say it:} \zh{我没有意见} — \emph{wǒ méiyǒu yìjiàn} — I've no objection.
  \item \emph{Recall:} say \emph{kànfǎ}, then \zh{意见}: a view, then a point you put forward
\end{itemize}

\begin{tcolorbox}[breakable,colback=teal!4,colframe=teal!35!black,title={Before you move on}]
What does \zh{意见} mean? (\textquotedblleft{}A suggestion; an objection\textquotedblright{}.) How do you say you have no objection? (\textbf{\zh{我没有意见}}.)
\end{tcolorbox}
\section[\texorpdfstring{\zh{意} (yì)}{意 (yì)}]{\zh{意} — the left of \zh{站} on top, then \zh{日} and \zh{心}}

\label{lesson:ZH-W280-yi}



\begin{quote}
The word first, then the character inside it.

\begin{itemize}
\raggedright
  \item \emph{Recall:} say \emph{yìjiàn} — \textbf{R1}, one lesson back, before it decays
\end{itemize}

\begin{quote}
\zh{站} \zh{日} \zh{心}
\end{quote}

\emph{Read it:} each one aloud, with nothing said first — \emph{zhàn}, \emph{rì}, \emph{xīn}
\end{quote}

\subsection*{Script — the left of \zh{站} on top, then \zh{日} and \zh{心}}
\hlblockfigure{\detokenize{figures/ZH-W280-yi-filmstrip.pdf}}{How \zh{意} is written, stroke by stroke}

\begin{quote}
\zh{意}
\end{quote}

\textbf{Three pieces, and you write every one of them already.}

\textbf{The left of \zh{站}, five strokes}, here on top and wider: the dot, the bar, a dot on the left, a left-falling stroke on the right, and a long horizontal. In \zh{站} that last stroke rises; here it lies flat and runs the whole width.

\textbf{\zh{日}, four strokes}, in the middle.

\textbf{\zh{心}, four strokes}, at the foot: the left dot, the curve along the foot that hooks up, the middle dot, the right dot.

5 + 4 + 4 = thirteen.

\subsection*{Your turn}
\begin{itemize}
\raggedright
  \item \emph{Write it:} \zh{意}, all thirteen, with the model in front of you
  \item \emph{Cover:} write it again from memory, and count the pieces, not the strokes
\end{itemize}

\begin{tcolorbox}[breakable,colback=teal!4,colframe=teal!35!black,title={Before you move on}]
How many strokes in \zh{意}? (\textbf{Thirteen}.) Which character you write has the same top? (\textbf{\zh{站}, where it stands on the left}.) What is at the foot? (\textbf{\zh{心}}.)

Source: \href{https://raw.githubusercontent.com/chanind/hanzi-writer-data/68d10a4b21150cae5e1ebbd223eed289cf32d90c/data/\%E6\%84\%8F.json}{Hanzi Writer Data \zh{意}, pinned stroke paths and medians}.
\end{tcolorbox}
\section[\texorpdfstring{\zh{我觉得}… (wǒ juéde …)}{我觉得… (wǒ juéde …)}]{\zh{我觉得}… — I think}

\label{lesson:ZH-C280-juede}



\begin{quote}
A word you have said for a long time, printed for the first time.

\begin{itemize}
\raggedright
  \item \emph{Read it:} \zh{意见}, aloud — \emph{yìjiàn}
  \item \emph{Recall:} say \emph{juéde}, to find; to think — \textbf{R3}, durable
  \item \emph{Recall:} say \emph{rènwéi}, to hold a view — \textbf{R3}
\end{itemize}
\end{quote}

\begin{grammarlens}[title={Grammar Lens: \zh{我觉得} and a whole sentence}]
\begin{quote}
\textbf{\zh{我觉得} and then a whole sentence gives your opinion: the sentence after it is what you think.}
\end{quote}

\begin{quote}
\textbf{\zh{我觉得他说得对}} \emph{wǒ juéde tā shuō de duì} — I think he's right. \textbf{\zh{我觉得今天太冷了}} \emph{wǒ juéde jīntiān tài lěng le} — I find it too cold today.
\end{quote}

Both characters of \emph{juéde}, \zh{觉} and \zh{得}, are ones you write, so \zh{我觉得} prints in characters now.

To ask, put \zh{你觉得} in front: \emph{nǐ juéde zěnmeyàng?}, what do you think?

\emph{rènwéi} says the same more formally, a view you hold and would defend; \zh{我觉得} is the everyday one.
\end{grammarlens}

\subsection*{Your turn}
Say these aloud:

\begin{itemize}
\raggedright
  \item \zh{我觉得他说得对} — \emph{wǒ juéde tā shuō de duì} — I think he's right.
  \item \zh{我觉得今天太冷了} — \emph{wǒ juéde jīntiān tài lěng le} — I find it too cold today.
  \item \emph{nǐ juéde zěnmeyàng} — What do you think?
  \item what you think of today's weather, with \zh{我觉得}
\end{itemize}

\begin{tcolorbox}[breakable,colback=teal!4,colframe=teal!35!black,title={Before you move on}]
What comes after \zh{我觉得}? (\textbf{A whole sentence: what you think}.) How do you ask what someone thinks? (\emph{nǐ juéde zěnmeyàng}.)
\end{tcolorbox}
\section[hǎoxiàng]{hǎoxiàng — as if; to look like}

\label{lesson:ZH-C280-haoxiang}



\begin{quote}
Before the new one: last lesson's opinion, and a word that sounds like this one's end.

\begin{itemize}
\raggedright
  \item \emph{Read it:} \zh{我觉得今天太冷了}, aloud — \emph{wǒ juéde jīntiān tài lěng le}
  \item \emph{Recall:} say \emph{xiàng}, to be like — \textbf{R3}, durable
  \item \emph{Recall:} say \zh{看来} — \emph{kànlai}, it seems — \textbf{R3}
\end{itemize}
\end{quote}

\subsection*{You'll want to know: hǎoxiàng}
\emph{hǎoxiàng} — \textquotedblleft{}as if; to look like\textquotedblright{}, a guess.

\begin{quote}
\emph{tā hǎoxiàng bìng le} — He looks as if he's ill.
\end{quote}

\begin{quote}
\emph{jīntiān hǎoxiàng yào xiàyǔ} — It looks like rain today.
\end{quote}

Its second character comes later; for now it is a word for the ear.

\begin{cousinweb}
\emph{hǎoxiàng} is built from \zh{好} and xiàng 'to be like', both inherited from Middle Chinese, \emph{xawX} and \emph{zjangX} (Guangyun, in Baxter's transcription). The dictionary glosses the first 'good, fine' and the second 'to be like'; 'as if' for the pair is the word's own sense.
\end{cousinweb}

\begin{grammarlens}[title={What you've built}]
\zh{我觉得} is your opinion. \zh{看来} draws a conclusion from what you see. \emph{hǎoxiàng} is softer still: a guess about how things look.

It goes before the verb, after the person: \emph{tā hǎoxiàng bìng le}.
\end{grammarlens}

\subsection*{Your turn}
\begin{itemize}
\raggedright
  \item \emph{Say it:} \emph{hǎoxiàng}
  \item \emph{Say it:} \emph{hǎoxiàng}, once more
  \item \emph{Say it:} \emph{tā hǎoxiàng bìng le} — He looks as if he's ill.
  \item \emph{Say it:} \emph{jīntiān hǎoxiàng yào xiàyǔ} — It looks like rain today.
  \item \emph{Recall:} say \emph{xiàng}, then \emph{hǎoxiàng}
\end{itemize}

\begin{tcolorbox}[breakable,colback=teal!4,colframe=teal!35!black,title={Before you move on}]
What does \emph{hǎoxiàng} mean? (\textquotedblleft{}As if; to look like\textquotedblright{}.) Is it sure? (\textbf{No: it is a guess}.)
\end{tcolorbox}
\section[\texorpdfstring{\zh{这么} (zhème)}{这么 (zhème)}]{\zh{这么} (zhème) — so; like this}

\label{lesson:ZH-C280-zheme}



\begin{quote}
Before the new one: last lesson's guess, and a word that points the other way.

\begin{itemize}
\raggedright
  \item \emph{Recall:} say \emph{hǎoxiàng}, as if — \textbf{R1}, one lesson back
  \item \emph{Read it:} \zh{那么}, aloud — \emph{nàme}
\end{itemize}
\end{quote}

\subsection*{You'll want to know: \zh{这么}}
\textbf{\zh{这么}} — \emph{zhème} — \textquotedblleft{}so; this much; like this\textquotedblright{}.

\begin{quote}
\textbf{\zh{今天这么冷}} \emph{jīntiān zhème lěng} — It's so cold today!
\end{quote}

\begin{quote}
\textbf{\zh{今天天气这么好}} \emph{jīntiān tiānqì zhème hǎo} — The weather's so good today!
\end{quote}

Both characters are ones you write.

\begin{cousinweb}
\textbf{\zh{这么}} (\emph{zhème}) is built from \zh{这} 'this' and the ending \zh{么} of \zh{什么}. No dictionary we use records an older form of \zh{这} in this sense, or of the ending, so its history is not settled.
\end{cousinweb}

\begin{grammarlens}[title={What you've built}]
\zh{那么} points at something over there, as in \zh{他没有你那么高}. \zh{这么} points at what is here and now: \zh{今天这么冷}.

Both go before an adjective: \zh{这么冷}, \zh{这么好}.
\end{grammarlens}

\subsection*{Your turn}
\begin{itemize}
\raggedright
  \item \emph{Say it:} \emph{zhème}
  \item \emph{Say it:} \emph{zhème}, once more
  \item \emph{Say it:} \zh{今天这么冷} — \emph{jīntiān zhème lěng} — It's so cold today!
  \item \emph{Say it:} \zh{今天天气这么好} — \emph{jīntiān tiānqì zhème hǎo} — The weather's so good today!
  \item \emph{Recall:} say \zh{那么}, then \zh{这么}: there, then here
\end{itemize}

\begin{tcolorbox}[breakable,colback=teal!4,colframe=teal!35!black,title={Before you move on}]
What does \zh{这么} mean? (\textquotedblleft{}So; this much\textquotedblright{}.) And \zh{那么}? (\textbf{So, of something over there}.)
\end{tcolorbox}
\section[Practice]{\zh{我觉得这个意见很好}}

\label{lesson:ZH-C280-practice}



\begin{quote}
The chapter closes here, with a plan and two people deciding on it.

\begin{itemize}
\raggedright
  \item \emph{Read it:} \zh{今天这么冷}, aloud — \emph{jīntiān zhème lěng}
  \item \emph{Recall:} say \emph{hǎoxiàng}, as if — \textbf{R2}
  \item \emph{Recall:} say \zh{我没有意见} — \emph{wǒ méiyǒu yìjiàn} — \textbf{R2}
\end{itemize}
\end{quote}

\subsection*{The exchange}
\begin{quote}
\textbf{A:} \zh{明天我们去山上} \zh{好不好}
\end{quote}

\begin{quote}
\emph{míngtiān wǒmen qù shān shàng, hǎo bu hǎo} — Shall we go up the mountain tomorrow?
\end{quote}

\begin{quote}
\textbf{B:} \zh{我觉得这个意见很好}
\end{quote}

\begin{quote}
\emph{wǒ juéde zhège yìjiàn hěn hǎo} — I think that's a good suggestion.
\end{quote}

\begin{quote}
\textbf{A:} \zh{你们有意见吗}
\end{quote}

\begin{quote}
\emph{nǐmen yǒu yìjiàn ma} — Any objections?
\end{quote}

\begin{quote}
\textbf{B:} \emph{wǒ méiyǒu yìjiàn, tā hǎoxiàng bù xiǎng qù} — None from me. He looks as if he doesn't want to go.
\end{quote}

\begin{quote}
\textbf{A:} \zh{为什么}
\end{quote}

\begin{quote}
\emph{wèishénme} — Why?
\end{quote}

\begin{quote}
\textbf{B:} \emph{wǒ juéde tā tài lèi le} — I think he's too tired.
\end{quote}

\begin{quote}
\textbf{A:} \zh{今天这么冷} \zh{明天会冷吗}
\end{quote}

\begin{quote}
\emph{jīntiān zhème lěng, míngtiān huì lěng ma} — It's so cold today. Will it be cold tomorrow?
\end{quote}

\begin{quote}
\textbf{B:} \zh{我觉得不会}
\end{quote}

\begin{quote}
\emph{wǒ juéde bú huì} — I don't think so.
\end{quote}

\zh{我觉得} answers on its own too: \zh{我觉得不会}, I don't think it will. The \zh{不} goes after \zh{我觉得}, with the verb it denies.

\begin{tcolorbox}[breakable,skin=enhanced,colback=orange!6,colframe=orange!45!black,boxrule=0.5pt,arc=1mm,left=6pt,right=6pt,top=4pt,bottom=4pt,fonttitle=\bfseries,title={Read this}]
\emph{Read it:} the exchange through once without stopping, then answer the question below
\end{tcolorbox}

\begin{grammarlens}[title={What you've built — giving an opinion}]
Four ways to say how you see it, from sure to unsure:

\begin{itemize}
\raggedright
  \item \textbf{\zh{我觉得他说得对}} \emph{wǒ juéde tā shuō de duì} — your opinion.
  \item \textbf{\zh{我没有意见}} \emph{wǒ méiyǒu yìjiàn} — no objection.
  \item \emph{tā hǎoxiàng bìng le} — a guess.
  \item \textbf{\zh{今天这么冷}} \emph{jīntiān zhème lěng} — so cold, here and now.
\end{itemize}

\textbf{Listen.} \emph{Nǐ juéde míngtiān huì xiàyǔ ma?}

\textbf{Speak.} Give your opinion of a plan for the weekend, and one objection.

\emph{Write it:} \zh{意见} — the left of \zh{站} on top, then \zh{日}, then \zh{心}
\end{grammarlens}

\subsection*{Your turn}
\begin{itemize}
\raggedright
  \item \emph{Say it:} \zh{我觉得这个意见很好} — \emph{wǒ juéde zhège yìjiàn hěn hǎo} — I think that's a good suggestion.
  \item \emph{Say it:} \emph{nǐ juéde zěnmeyàng} — What do you think?
  \item \emph{Say it:} the whole exchange, both parts
  \item \emph{Write it:} \zh{我觉得今天这么冷}
\end{itemize}

\begin{tcolorbox}[breakable,colback=teal!4,colframe=teal!35!black,title={Before you move on}]
Say you think he's right. (\textbf{\zh{我觉得他说得对}}.) Say you have no objection. (\textbf{\zh{我没有意见}}.) Say he looks ill. (\emph{tā hǎoxiàng bìng le}.)
\end{tcolorbox}
